\section{Reading a tactic failure}\label{sec:04-tactics:03-reading-failures:1}

Every tactic either makes progress or produces an error, and the error message is almost always telling you exactly what to do next. Treat it as a diagnostic, not a dead end.

\begin{itemize}
\tightlist
\item
  \href{https://lean-lang.org/doc/reference/latest/Tactic-Proofs/Tactic-Reference/}{\texttt{rw [h]}} where \texttt{h : a = b} but the goal does not literally contain \texttt{a} fails with \textbf{``motive is not type correct''} or ``did not find instance of the pattern.'' This means the exact syntactic term \texttt{a} being rewritten does not occur in the goal as written. Fix: use \texttt{\#check} (or hover in the editor) on the goal's actual statement. It is common for a definition to be unfolded differently than expected, so the term to be rewritten is hiding behind a \texttt{def} that needs \href{https://lean-lang.org/doc/reference/latest/Tactic-Proofs/Tactic-Reference/}{\texttt{unfold}} or \href{https://lean-lang.org/doc/reference/latest/Tactic-Proofs/Tactic-Reference/}{\texttt{show}} first.
\item
  \href{https://lean-lang.org/doc/reference/latest/Tactic-Proofs/Tactic-Reference/}{\texttt{exact e}} where \texttt{e}'s type does not match the goal fails with a \textbf{type mismatch}, printing both the expected and the actual type side by side. The difference between them almost always points to one wrong argument or a missing \texttt{.symm}.
\item
  \href{https://lean-lang.org/doc/reference/latest/Tactic-Proofs/Tactic-Reference/}{\texttt{apply f}} where \texttt{f}'s conclusion does not unify with the goal fails in a similar way, but \emph{before} spending effort on \texttt{apply}, \texttt{\#check f} shows its full type. This reveals exactly what subgoals \texttt{apply} would leave, before committing to the tactic.
\item
  A tactic that produces \textbf{no error but does not close the goal} is not a failure. The resulting goal state should be checked (in the editor, or with \href{https://lean-lang.org/doc/reference/latest/Tactic-Proofs/Tactic-Reference/}{\texttt{sorry}} in place of the rest of the proof, which type-checks with a warning and permits inspection of exactly what remains).
\end{itemize}

\texttt{sorry} deserves a special mention: it is a placeholder that closes any goal instantly but marks the theorem as unproved. (Lean prints a warning, and any downstream proof depending on it inherits this unproved status.) It should be used constantly while exploring. Writing the skeleton of a multi-step proof with \texttt{sorry} at each unfinished branch, checking that the overall \emph{shape} type-checks, and then filling in the branches one at a time, is normal practice, not a hack.
